\section{Проверка}
\label{sec:verification}

Проверката има три нива. Тестовете на ниво функция проверяват частите поотделно. Тестът
на TLS сесията провежда истински ръкостискания между клиент на OpenSSL и сесията на
\code{octet} в един процес, без сокети. Интеграционният тест стартира изпълнимия файл и
говори с него през истински сокети, така, както би говорил реален клиент. Всички тестове
се стартират с \code{ctest} и минават на двете проверени платформи
\tabref{tab:platforms}. Пълните записи от изпълненията са в хранилището, в папката
\code{evidence}.

\begin{table}[H]
\centering
\caption{Платформи, на които тестовете минават}
\label{tab:platforms}
\begin{tabularx}{\textwidth}{@{}lYY@{}}
\toprule
 & Windows & Linux \\
\midrule
Система & Windows 11, 10.0.26200 & Ubuntu 24.04 във WSL2, ядро 6.6 \\
Компилатор & MSVC 19.51 & GCC 14.2 \\
OpenSSL & 3.6.4 (vcpkg) & 3.0.13 (от дистрибуцията) \\
Предупреждения при компилиране & няма при \code{/W4 /permissive-} & няма при
  \code{-Wall -Wextra -Wpedantic -Wconversion -Wshadow} \\
\bottomrule
\end{tabularx}
\end{table}

\subsection{Тестове на ниво функция}

Тестовете на ниво функция са три програми с общо \checksUnitTotal{} проверки, без
външна рамка за тестване.

Тестът на разпознаването (\checksClassify{} проверки) минава през всичките 256 възможни
стойности на първия октет и проверява, че точно една от тях води до TLS и точно 26 до
HTTP. Отделно проверява всички стандартни методи на HTTP, началната последователност на
HTTP/2, трите TLS записа, които не могат да отворят връзка, стария формат на SSL 2.0,
празен вход, малки букви и водещ интервал. Проверява и че разпознаването гледа само
първия октет: 0x16 в средата на заявка не променя решението.

Тестът на връщането на октети (\checksReplay{} проверки) проверява, че прочетените при
разпознаването октети излизат първи, че източникът не се чете, докато те не свършат, и
че TLS запис, пристигнал по един октет, излиза цял и в правилния ред.

Тестът на анализатора на HTTP (\checksHttp{} проверки) проверява валидни заглавни части
и десет начина на неправилна: малки букви в метода, цел без наклонена черта, версия
HTTP/2.0, двоен интервал, излишна дума, управляващ знак, твърде дълъг метод, сгъване на
ред, интервал преди двоеточие и поле без двоеточие. Проверява и заглавна част, която
надхвърля границата, подадена и по един октет, и на големи части.

\subsection{TLS сесия в паметта}

Тестът на TLS (\checksTls{} проверки) провежда истински ръкостискания без мрежа и без
втора нишка. Клиентът е обикновен клиент на OpenSSL с проверка на сертификата и на
името. Хитростта е във функцията, от която сесията на сървъра чете: когато сървърът има
нужда от октети и няма такива, тя изпълнява една стъпка от клиента. Клиентът обработва
каквото сървърът е изпратил и произвежда следващата си група съобщения. Двете страни се
редуват в една нишка.

Първият октет на ClientHello се изважда и се подава обратно чрез \code{ReplayReader},
точно както прави сървърът след разпознаването. Проверени са:
\begin{itemize}
\item ръкостискане, при което сървърът получава ClientHello по един октет на четене;
  договорена е версия TLS 1.3 и протокол \code{http/1.1}, а клиентът потвърждава веригата
  до тестовия орган;
\item обмен на данни в двете посоки след ръкостискането;
\item клиент, който предлага само \code{h2}, получава предупреждение 120;
\item клиент без ALPN и сървър в нестрог режим се приемат без избран протокол;
\item клиент до TLS 1.2 се приема при минимум 1.2 и се отказва при минимум 1.3;
\item клиент, който очаква друго име, отхвърля сертификата с грешка за несъвпадащо име.
\end{itemize}

\subsection{Интеграционен тест}

Интеграционният тест (\checksIntegration{} проверки) стартира изпълнимия файл в шест
конфигурации и говори с него с модула \code{ssl} на Python за TLS и с обикновени сокети
за всичко останало. Проверките са групирани по конфигурации в
таблица~\ref{tab:integration}.

\begin{table}[H]
\centering
\caption{Конфигурации и проверки на интеграционния тест}
\label{tab:integration}
\begin{tabularx}{\textwidth}{@{}L{4.6cm}Y@{}}
\toprule
Конфигурация & Какво се проверява \\
\midrule
Един сокет, разпознаване, сертификат &
  TLS 1.3 с \code{http/1.1}; некриптиран HTTP на същия порт; без отговор на двоични данни,
  на малки букви и на предупреждение като първи запис; 501 за непознат метод; 400 за
  неправилен ред; 431 за огромна заглавна част; 0x16 с безсмислени данни не получава
  некриптиран отговор; отказ на клиент само с \code{h2} и точните октети на
  предупреждението; клиент без ALPN се обслужва; мълчащ клиент не пречи на другите и се
  затваря при срока; клиент, който спира след 0x16, се затваря при същия срок; записът
  на сървъра съдържа всички очаквани изходи \\
ClientHello по един октет & ръкостискането завършва; разпознаването е станало по част
  от ClientHello, а останалото е пристигало дълго след това \\
Без сертификат & TLS клиентът се затваря без некриптиран отговор; HTTP се обслужва \\
Минимум TLS 1.3 & клиент до TLS 1.2 се отказва; клиент с TLS 1.3 се обслужва \\
Отделни слушатели & слушател само за TLS не отговаря на HTTP; слушател само за HTTP не
  завършва ръкостискане \\
Граница от две връзки & третата връзка се затваря веднага \\
\bottomrule
\end{tabularx}
\end{table}

Проверката на отказа по ALPN заслужава отделно внимание. Python на Linux е свързан с
по-стара версия на OpenSSL, която не разпознава името на това предупреждение и съобщава
„unknown error“. Затова тестът не разчита на текста на грешката. Той изпраща ClientHello
през обикновен сокет и търси в отговора точните седем октета:

\begin{center}
\code{15 03 03 00 02 02 78}
\end{center}

Това е запис с ContentType 21 (предупреждение), версия 0x0303, дължина 2, ниво 2
(фатално) и описание 120 (no\_application\_protocol). Записът идва некриптиран, защото
сървърът отказва още при обработката на ClientHello, преди ServerHello. Тази проверка
минава еднакво и на двете платформи.

Тестът с ClientHello по един октет също е показателен. Клиентът изпраща всеки октет в
отделен TCP сегмент с пауза от една милисекунда. Сървърът записва колко октета е донесло
първото четене. Броят зависи от планирането на нишките, затова тестът проверява само, че
е по-малък от сто, тоест че разпознаването е станало по малка част от ClientHello, и че
до готовността на транспорта са минали над 300 милисекунди. Ръкостискането завършва
успешно.

\subsection{Какво остава непроверено}

Конфигурацията за macOS съществува, но не е компилирана и изпълнена. Тестовете не
включват автоматизирано подаване на случайни входни данни. Тестовете на Linux са
изпълнени във WSL2, тоест във виртуална машина, а не на физически компютър с Linux.
